
The \texttt{agentical} Rust crate has no dependencies beyond the standard library.
It defines synchronous, executor-agnostic traits for the agent, the oracle, and the journal,
shown in Listing~\ref{lst:traits}, along with interfaces for channels, memory, and supervision
intended for future multi-agent work (Section~\ref{sec:limitations}).

\begin{lstlisting}[caption={The agent, oracle, and journal traits (doc comments omitted).},label={lst:traits}]
pub trait Agent {
    type State;
    type Out;
    fn step(&mut self, state: &mut Self::State) -> Step<Self::Out>;
}

pub trait Oracle {
    fn sample(&self, prompt: Prompt) -> Sampled;
    fn measure_hint(&self) -> Option<Dist> { None }
}

pub trait Journal {
    fn append(&mut self, agent: &AgentId, event: Event) -> Cursor;
    fn replay_next(&self, agent: &AgentId) -> Outcome;
    fn checkpoint(&self, agent: &AgentId) -> Cursor;
}
\end{lstlisting}

The journal is position-indexed: the cursor is the key, as in Section~\ref{sec:modes}, so there
is no key-to-value map and no collision assumption.
A runtime module implements the semantics for a single agent:
\begin{itemize}[noitemsep]
  \item \texttt{InMemoryJournal}, with \texttt{append}, \texttt{replay\_next}, \texttt{rewind},
    and \texttt{truncate} (to simulate a crash);
  \item \texttt{SeededOracle}, a categorical sampler driven by a SplitMix64 generator, which
    stands in for an LLM;
  \item a driver that runs an agent in record, replay, or resume mode following
    \textsc{[O-rec]}, \textsc{[O-rep]}, and \textsc{[O-resume]}.
\end{itemize}

\paragraph{Side effects.}
The prototype also journals tool side effects, outside the formal results, using a
write-ahead discipline.
Before performing an effect it appends an \emph{intent} record carrying the effect's argument;
after the effect succeeds it appends a \emph{committed} record carrying the result.
Replay returns the journaled result and does not perform the effect again.
A crash between the two records leaves a dangling intent, which recovery detects and resolves
either by performing the effect again with the journaled argument, which is safe for an
idempotent effect of a single agent, or by recording a compensation.
Replaying a compensated effect stops the run without performing anything.
Section~\ref{sec:eval-e2} tests this behaviour at every crash point; formalizing it is future
work.

\paragraph{What is not implemented.}
The channel, memory, and supervisor interfaces have no implementation, so the prototype runs a
single agent.
The journal lives in memory only; there is no persistent backend.
